\section{Unidades}

Use preferencialmente unidades do Sistema Internacional (SI). Outras unidades
podem ser utilizadas como unidades secundárias (entre parênteses). Isto se aplica
em armazenamento de dados. Escreva,
por exemplo,  ``$15\,\mathrm{Gb}/\mathrm{cm}^2$ ($100\,\mathrm{Gb}/\mathrm{in}^2$)''. 
Uma exceção é quando unidades inglesas são usadas na identificação de algum item
comercial, como ``disco de 3.5 polegadas''. Evite misturar unidades SI com outras, 
como por exemplo corrente elétrica em amperes e campo magnético em oersteds. 
Isso frequentemente causa confusão, pois as equações não batem dimensionalmente. 
Se você realmente precisar usar unidades mistas, deixe claro as unidades para cada
termo em uma equação. A unidade SI para força do campo magnético 
$\mathbf{H}$ é $\mathrm{A}/\mathrm{m}$. No entanto, se você quiser utilizar unidades
de  $\mathrm{T}$, 
refira-se à densidade de fluxo magnético $\mathbf{B}$ ou
força do campo magnético simbolizada por $\mu_0\,\mathbf{H}$. 
Use um ponto centralizado para separar unidades compostas, p. ex.,  
``$\mathrm{A} \cdot \mathrm{m}^2$''.

%Use SI as primary units. Other units may be used as secondary units
%(in parentheses). This applies to papers in data storage. For example,
%write ``$15\,\mathrm{Gb}/\mathrm{cm}^2$ ($100\,\mathrm{Gb}/\mathrm{in}^2$)''. 
%An exception is when
%English units are used as identifiers in trade, such as ``3.5 in
%disk drive''. Avoid combining SI and other units, such as current in
%amperes and magnetic field in oersteds. This often leads to confusion
%because equations do not balance dimensionally. If you must use mixed
%units, clearly state the units for each quantity in an equation.  The
%SI unit for magnetic field strength $\mathbf{H}$ is $\mathrm{A}/\mathrm{m}$. However, if you wish to
%use units of $\mathrm{T}$, either refer to magnetic flux density $\mathbf{B}$ or
%magnetic field strength symbolized as $\mu_0\,\mathbf{H}$. Use the center dot to
%separate compound units, e.g., ``$\mathrm{A} \cdot \mathrm{m}^2$''.

\section{Dicas Úteis}

\subsection{Figuras e Tabelas}

Legendas dos eixos das figuras são frequentemente uma fonte de confusões. Use
palavras ao invés de símbolos. Por exemplo, escreva a quantidade ``Magnetização'',
ou ``Magnetização M'', não apenas ``M''. Coloque unidades entre parênteses. Não marque
os eixos apenas com unidades ou quantidades. Por exemplo, escreva 
``Temperatura ($\mathrm{K}$)'', não ``$\mbox{Temperature}/\mathrm{K}$''.

Multiplicandos podem ser especialmente confusos. Escreva  ``Magnetização 
($\mathrm{kA}/\mathrm{m}$)'' ou ``Magnetização ($10^3 \mathrm{A}/\mathrm{m}$)''. 
Não escreva 
``Magnetização $(\mathrm{A}/\mathrm{m}) \times 1000$'' pois o leitor não saberá
se a legenda do eixo significa 
$16000\,\mathrm{A}/\mathrm{m}$ ou $0.016\,\mathrm{A}/\mathrm{m}$.

%Figure axis labels are often a source of confusion. Use words rather
%than symbols. As an example, write the quantity ``Magnetization'', or
%``Magnetization M'', not just ``M''. Put units in parentheses. Do not
%label axes only with units.  For example, write ``Magnetization
%($\mathrm{A}/\mathrm{m}$)'' or ``Magnetization ($\mathrm{A} \mathrm{m}^{-1}$)'', not just
% ``$\mathrm{A}/\mathrm{m}$''. Do not
%label axes with a ratio of quantities and units. For example, write
%``Temperature ($\mathrm{K}$)'', not ``$\mbox{Temperature}/\mathrm{K}$''.
%
%Multipliers can be especially confusing. Write ``Magnetization
%($\mathrm{kA}/\mathrm{m}$)'' or ``Magnetization ($10^3 \mathrm{A}/\mathrm{m}$)''. Do not write
%``Magnetization $(\mathrm{A}/\mathrm{m}) \times 1000$'' because the reader would not know
%whether the axis label means $16000\,\mathrm{A}/\mathrm{m}$ or $0.016\,\mathrm{A}/\mathrm{m}$.

\subsection{Referências}

Use referências no estilo Harvard (veja no final deste documento). Com
\LaTeX, você pode processar uma base externa de dados bibliográficos 
usando \texttt{bibtex},\footnote{Nesse caso, você também vai precisar do arquivo \texttt{ifacconf.bst}, que é parte do pacote \texttt{ifacconf\_latex.zip}} ou então inserir
as referências diretamente na seção correspondente. Evite as notas de rodapé.
Por favor, note que o estilo preferido de referências é o que se encontra no
final deste documento. Artigos que não foram publicados devem ser citados 
como ``unpublished''. Use maiúscula apenas na primeira palavra do título
do artigo, exceto para nomes próprios, acrônimos e símbolos. Todos os autores
devem ser explicitados nos artigos. Note que você pode citar com \texttt{\textbackslash cite\{label\}} quando a citação fizer parte do texto ou com \texttt{\textbackslash citep\{label\}}.

%Use Harvard style references (see at the end of this document). With
%\LaTeX, you can process an external bibliography database 
%using \texttt{bibtex},\footnote{In this case you will also need the \texttt{ifacconf.bst}
%file, which is part of the \texttt{ifaconf} package.}
%or insert it directly into the reference section. Footnotes should be avoided as
%far as possible.  Please note that the references at the end of this
%document are in the preferred referencing style. Papers that have not
%been published should be cited as ``unpublished''.  Capitalize only the
%first word in a paper title, except for proper nouns and element
%symbols.

\subsection{Abreviações e Acrônimos}

Defina as abreviações e os acrônimos na primeira ocorrência do texto do artigo, 
mesmo se foram definidos no abstract. Abreviações conhecidas, como SI ou SBA, não
precisam ser definidas. Abreviações que incorporam pontos não devem conter espaços
(escreva ``C.N.R.S.'', não ``C. N. R. S.''). Não use abreviações no título, exceto
se forem inevitáveis. 
%Define abbreviations and acronyms the first time they are used in the
%text, even after they have already been defined in the
%abstract. Abbreviations such as IFAC, SI, ac, and dc do not have to be
%defined. Abbreviations that incorporate periods should not have
%spaces: write ``C.N.R.S.'', not ``C. N. R. S.'' Do not use abbreviations
%in the title unless they are unavoidable (for example, ``IFAC'' in the
%title of this article).

\subsection{Equações}

Numere as equações de forma consecutiva, com números entre parênteses
encostados no lado direito da
margem, como em~\eqref{eq:sample}. Para tornar suas equações mais compactas,
você pode usar $/$, a função $\exp$, ou expoentes apropriados. Use parênteses para
evitar ambiguidades. Pontue as equações quando estas forem parte do texto, como em
%Number equations consecutively with equation numbers in parentheses
%flush with the right margin, as in (\ref{eq:sample}).  To make your equations more
%compact, you may use the solidus ($/$), the $\exp$ function, or
%appropriate exponents. Use parentheses to avoid ambiguities in
%denominators. Punctuate equations when they are part of a sentence, as
%in

\begin{multline} 
\int_0^{r_2}  F (r, \varphi ) dr d\varphi =  [\sigma r_2 / (2 \mu_0 )] \\
\times   \int_0^{\infty} \exp(-\lambda |z_j - z_i |) \lambda^{-1} J_1 (\lambda  r_2 ) J_0 (\lambda r_i ) d\lambda 
\label{eq:sample2}
\end{multline}

Certifique-se que os símbolos em sua equação tenham sido previamente definidos (ou defina-se imediatamente depois). Símbolos devem aparecer em itálico ($T$ pode
se referir à temperatura, mas T é a unidade tesla). 
Refira-se a ``(\ref{eq:sample})'', não ``Eq. (\ref{eq:sample})'' ou ``equação
(\ref{eq:sample})'', exceto no começo de uma frase: ``Equação~(\ref{eq:sample}) é \ldots''. 
%might refer to temperature, but T is the unit tesla). 
%Be sure that the symbols in your equation have been defined before the
%equation appears or immediately following. Italicize symbols ($T$
%might refer to temperature, but T is the unit tesla). Refer to
%``(\ref{eq:sample})'', not ``Eq. (\ref{eq:sample})'' or ``equation
%(\ref{eq:sample})'', except at the beginning of a sentence: ``Equation
%(\ref{eq:sample}) is \ldots''.

\subsection{Outras Recomendações}

Deixe um espaço após pontos e vírgulas. 
Uma afirmação entre parênteses no final de uma sentença é pontuada exteriormente
(assim). (Uma frase entre parênteses é pontuada internamente.) Evite contrações
(por exemplo, ``prá'' chegar ao resultado --- use ``para'').
